\chapter{결론 및 향후 연구}
\label{ch:conclusion}

\section{결론}

본 논문은 분 단위 공개 텔레메트리라는 제약 아래에서 리그 오브 레전드 교전 결과 예측 문제를 정식화하고, 이를 위한 완전한 파이프라인을 제안·평가하였다. 킬 조건부 교전 국소화 알고리즘은 사람의 주석 없이 206{,}442 경기로부터 약 100 만 건의 교전을 결정론적으로 찾아내며, 교환가치 라벨은 킬·오브젝트·현상금이 섞인 교전의 순이득을 이진 결과로 요약한다. 무누출 계약을 명시한 다중 모달 표현 위에서 25 종 이상의 모델을 동일 조건으로 비교한 결과, 공학적 테이블 요약과 그래디언트 부스팅의 조합이 시험 AUC $.675$로 가장 우수하였고, 시퀀스·그래프·사건 기반 신경망 표현은 $.569$--$.581$에 머물렀다. 정합 입력 비교는 이 격차의 절반이 학습기 계열에서, 나머지 절반이 표현에서 비롯됨을 보였다. 예측 가능성은 후반 국면과 일방적 골드 격차에서 $.8$에 이르렀으나 대등한 교전에서는 우연 수준에 근접하였는데, 이는 관측의 정보 입도가 부과하는 상한으로 해석된다. 연구 질문에 대한 답은 다음과 같이 요약된다.

\begin{description}[leftmargin=3.2em,style=nextline]
  \item[RQ1] 킬 조건부 규칙만으로 재현 가능한 교전 말뭉치를 구축할 수 있다. 사람 주석과의 일치도는 두 트랙의 검증 프로토콜로 계량되며 결과는 진행 중이다.
  \item[RQ2] 공학적 테이블 요약이 가장 많은 신호를 드러낸다. 시퀀스·그래프·사건 표현의 추가는 신호를 더 드러내지 못했다.
  \item[RQ3] 입력을 고정했을 때 그래디언트 부스팅이 MLP보다 $0.049$ 높다.
  \item[RQ4] 예측 가능성은 국면·골드 격차에 따라 단조 증가하며, 대등한 교전의 낮은 예측 가능성은 정보 입도 한계로 설명된다.
  \item[RQ5] 일곱 처치의 효과는 절제 프로토콜로 계량된다. \todo{최종 실행 결과 반영.}
\end{description}

\section{향후 연구}

\begin{itemize}[leftmargin=1.5em]
  \item \textbf{검출기 검증의 완료.} 2026 리플레이 골드셋과 2025 패치에 대한 사람 주석을 마치고 정밀도·재현율·$\kappa$를 보고한다.
  \item \textbf{정보 입도의 직접 측정.} 리플레이로부터 초 단위 관측을 얻어 같은 교전에 대해 해상도를 단계적으로 낮추며 AUC의 변화를 측정하면, 본 논문이 추론한 정보 입도 한계를 실험적으로 확인할 수 있다.
  \item \textbf{챔피언 능력 부호화.} 스킬 세트, 쿨다운, 아이템 효과를 구조화된 특징으로 도입한다.
  \item \textbf{교차 지역·프로 경기 전이.} KR 외 서버와 프로 경기에 대한 전이 학습을 평가한다.
  \item \textbf{인과 분석.} 예측을 넘어 교전 결과의 인과 요인을 식별한다.
  \item \textbf{실시간 코칭.} 추론 지연을 줄여 경기 중 의사결정 보조에 적용한다.
\end{itemize}
